\section{Exact determinant and renewal equations}
\label{sec:exact}

Let $A_N$ be the adjacency matrix of $G_N$ and define
\[
 D(z)=\det(I-zA),\qquad
 C_v(z)=\det(I-zA^{(v)}).
\]
The empty determinant is one when $A$ is $1\times1$.
The graph enlargement obtained by adding the new $N$-edge chain is defined
for every $N\ge1$, independently of whether an old length-$N$ return walk
has been selected.  The determinant statements below therefore hold for
every $N\ge1$; the split factor of \cref{thm:construction} is asserted only
for $N\ge N_0$, when the required image walk exists.

\begin{proposition}[Path-elimination identity]
\label{prop:determinant}
For every $N\ge1$,
\begin{equation}
 \det(I-zA_N)=D(z)-z^NC_v(z).
 \label{eq:determinant}
\end{equation}
\end{proposition}

\begin{proof}
Assume first $N\ge2$ and order the old vertices before the $N-1$ new
vertices along the chain.  Let $Q$ be the nilpotent adjacency matrix of the
directed chain on the new vertices and put $R=I-zQ$.  Then $\det R=1$ and
\[
 (R^{-1})_{1,N-1}=z^{N-2}.
\]
Writing coordinate vectors for the relevant endpoints gives
\[
 I-zA_N=
 \begin{pmatrix}
 I-zA & -z e_ve_1^T\\
 -z e_{N-1}e_v^T & R
 \end{pmatrix}.
\]
The Schur complement of $R$ is
\[
 I-zA-z^2(R^{-1})_{1,N-1}e_ve_v^T
 =I-zA-z^Ne_ve_v^T.
\]
The determinant is linear in its $(v,v)$ entry, whose cofactor is $C_v(z)$.
Since $\det R=1$, \cref{eq:determinant} follows.  When $N=1$, the operation
adds a parallel loop at $v$, and the same identity is the corresponding
one-entry determinant update.
\end{proof}

Let $f_n$ count old length-$n$ paths from $v$ to $v$ that avoid $v$ at
times $1,\ldots,n-1$, and put
\[
 F_v(z)=\sum_{n\ge1}f_nz^n.
\]
Every based loop decomposes uniquely into first-return loops.  Hence, as
formal power series,
\begin{align}
 L_v(z)
 &:=\sum_{n\ge0}(A^n)_{vv}z^n
 =[(I-zA)^{-1}]_{vv} \notag\\
 &=\frac{C_v(z)}{D(z)}
 =1+F_v(z)+F_v(z)^2+\cdots
 =\frac1{1-F_v(z)}.
 \label{eq:renewal}
\end{align}
The determinant ratio is a standard renewal identity; compare
\cite{Gomez2010,BlockEtAl1980}.  In particular,
\begin{equation}
 F_v(z)=1-\frac{D(z)}{C_v(z)}.
 \label{eq:F-determinant}
\end{equation}

\begin{proposition}[Entropy-radius equation]
\label{prop:radius}
Let $r=\rho(A)^{-1}$ and $r_N=\rho(A_N)^{-1}$.  Then $F_v$ is analytic on a
neighborhood of $[0,r]$, $F_v(r)=1$, and $r_N$ is the unique point in
$(0,r)$ satisfying
\begin{equation}
 F_v(r_N)+r_N^N=1.
 \label{eq:radius}
\end{equation}
\end{proposition}

\begin{proof}
Since $A$ is irreducible, the proper principal submatrix $A^{(v)}$ has
$\rho(A^{(v)})<\rho(A)$.  Thus $C_v$ has no zero at $r$, the first-return
series is analytic beyond $r$, and \cref{eq:F-determinant} with $D(r)=0$
gives $F_v(r)=1$.

All internal vertices of the added chain are new and different from $v$.
The chain is therefore one new first-return path and changes the series to
$F_v(z)+z^N$.  Its old image may revisit $v$; that does not alter this count.
The new series is continuous and strictly increasing on $[0,r]$, takes value
zero at zero, and exceeds one at $r$.  It consequently has a unique root
$s_N\in(0,r)$ of \cref{eq:radius}.

\Cref{eq:determinant,eq:F-determinant} show that $s_N$ is a
positive zero of $\det(I-zA_N)$.  For $N\ge2$, the old matrix is a proper
principal block of the irreducible $A_N$, so strict Perron monotonicity gives
$\rho(A_N)>\rho(A)$.  For $N=1$, one has
$A_1=A+e_ve_v^T$, so $A_1\ge A$ and $A_1\ne A$; the strict comparison
theorem for irreducible nonnegative matrices gives the same inequality.  Thus the
reciprocal Perron root lies in $(0,r)$ and is the smallest positive
determinant zero.  Hence $s_N=r_N$.
\end{proof}

\begin{remark}
\Cref{eq:determinant} and \cref{eq:radius} describe the same perturbation,
but they protect different errors.  The Schur complement fixes the minus sign
and exponent; the renewal equation explains why the old image walk need not
itself be first return.
\end{remark}
